
\documentclass[12pt,a4paper]{article}

% =========================================================
% PAQUETES
% =========================================================
\usepackage[utf8]{inputenc}
\usepackage[T1]{fontenc}
\usepackage[spanish,es-nodecimaldot]{babel}
\usepackage{lmodern}
\usepackage{geometry}
\usepackage{graphicx}
\usepackage{booktabs}
\usepackage{array}
\usepackage{longtable}
\usepackage{float}
\usepackage{amsmath,amssymb}
\usepackage{enumitem}
\usepackage{xcolor}
\usepackage{titlesec}
\usepackage{fancyhdr}
\usepackage{hyperref}
\usepackage{setspace}
\usepackage{lastpage}
\usepackage{caption}
\usepackage{subcaption}
\usepackage{tikz}
\usepackage{tabularx}
\usepackage{listings}
\usepackage[htt]{hyphenat}

% =========================================================
% CONFIGURACIÓN GENERAL
% =========================================================
\geometry{
    top=2.5cm,
    bottom=2.5cm,
    left=2.7cm,
    right=2.7cm
}

\setstretch{1.15}
\setlength{\parindent}{0pt}
\setlength{\parskip}{0.6em}

% Color institucional aproximado para títulos.
% Puede modificarse si se desea emplear otro tono.
\definecolor{UBUblue}{RGB}{0,82,147}
\definecolor{lightgray}{RGB}{245,245,245}
\definecolor{UBUlight}{RGB}{235,243,249}
\definecolor{UBUdark}{RGB}{34,45,58}

% Carpeta donde se encuentra el logotipo institucional

\hypersetup{
    colorlinks=true,
    linkcolor=UBUblue,
    urlcolor=UBUblue,
    citecolor=UBUblue,
    pdftitle={Práctica I - Refactorización de una planta industrial mediante POO},
    pdfauthor={Universidad de Burgos}
}

% =========================================================
% DATOS DEL TRABAJO
% =========================================================

\newcommand{\universidad}{Universidad de Burgos}
\newcommand{\grado}{Grado en Tecnologías Digitales para la Empresa}
\newcommand{\asignatura}{Programación Modular y Orientación a Objetos}
\newcommand{\titulotrabajo}{Práctica I: Refactorización de una planta industrial}

\newcommand{\nombrealumno}{Alejandro Esteban Gómez}
\newcommand{\dni}{71565225N}
\newcommand{\correo}{aeg1014@alu.ubu.es}
\newcommand{\grupo}{Grupo 1}
\newcommand{\profesor}{Jaime Andres Rincon Arango}
\newcommand{\cursoacademico}{2026-2027}
\newcommand{\fechaentrega}{2 Noviembre 2026}

% Enlace al repositorio GitHub con el código completo (OBLIGATORIO)
\newcommand{\repositorio}{https://github.com/aeg1014/PR1_POO.git}


% =========================================================
% FORMATO DE TÍTULOS
% =========================================================

\titleformat{\section}
  {\Large\bfseries\color{UBUblue}}
  {\thesection.}{0.6em}{}

\titleformat{\subsection}
  {\large\bfseries\color{UBUblue}}
  {\thesubsection.}{0.6em}{}

\titleformat{\subsubsection}
  {\normalsize\bfseries}
  {\thesubsubsection.}{0.6em}{}

% =========================================================
% CABECERA Y PIE
% =========================================================

\pagestyle{fancy}
\fancyhf{}

\fancyhead[L]{\small \asignatura}
\fancyhead[R]{\small Práctica I}

\fancyfoot[L]{\small \universidad}
\fancyfoot[C]{\small \nombrealumno}
\fancyfoot[R]{\small Página \thepage\ de \pageref{LastPage}}

\setlength{\headheight}{15pt}
\renewcommand{\headrulewidth}{0.4pt}
\renewcommand{\footrulewidth}{0.4pt}


% =========================================================
% ESTILO DE CÓDIGO (fragmentos de Python)
% =========================================================
\lstdefinestyle{python}{
    language=Python,
    basicstyle=\ttfamily\footnotesize,
    keywordstyle=\color{UBUblue}\bfseries,
    commentstyle=\color{gray}\itshape,
    stringstyle=\color{UBUdark!70!red},
    backgroundcolor=\color{lightgray},
    frame=single,
    rulecolor=\color{UBUlight!60!UBUblue},
    numbers=left,
    numberstyle=\tiny\color{gray},
    breaklines=true,
    showstringspaces=false,
    tabsize=4,
    xleftmargin=1.5em,
    framexleftmargin=1.5em,
    literate={á}{{\'a}}1 {é}{{\'e}}1 {í}{{\'i}}1 {ó}{{\'o}}1 {ú}{{\'u}}1
             {Á}{{\'A}}1 {É}{{\'E}}1 {Í}{{\'I}}1 {Ó}{{\'O}}1 {Ú}{{\'U}}1
             {ñ}{{\~n}}1 {Ñ}{{\~N}}1 {ü}{{\"u}}1 {¿}{{?`}}1 {¡}{{!`}}1
}
\lstset{style=python}

% Permite partir identificadores largos (p. ej. \texttt{objeto\_detectado}) tras el guion bajo
\let\guionbajo\_
\renewcommand{\_}{\guionbajo\allowbreak}
\setlength{\emergencystretch}{3em}
\renewcommand{\lstlistingname}{Código}

\graphicspath{{logo/}{figuras/}}

% =========================================================
% DOCUMENTO
% =========================================================

\begin{document}

% =========================================================
% PORTADA
% =========================================================

\begin{titlepage}
    \thispagestyle{empty}

    % Elementos gráficos de fondo
    \begin{tikzpicture}[remember picture,overlay]
        % Banda institucional izquierda
        \fill[UBUblue]
            (current page.north west)
            rectangle
            ([xshift=1.15cm]current page.south west);

        % Línea superior de acento
        \fill[UBUblue]
            ([xshift=1.15cm,yshift=-0.8cm]current page.north west)
            rectangle
            ([yshift=-0.95cm]current page.north east);

        % Bloque inferior suave
        \fill[UBUlight]
            ([xshift=1.15cm,yshift=4.0cm]current page.south west)
            rectangle
            (current page.south east);
    \end{tikzpicture}

    \hspace*{0.8cm}
    \begin{minipage}[t]{0.89\textwidth}

        \vspace*{0.2cm}

        % -------------------------------------------------
        % LOGO UBU
        % El archivo debe estar en: logo/logo_ubu.png
        % -------------------------------------------------
        \IfFileExists{logo/logo_ubu.png}{%
            \includegraphics[width=0.64\textwidth]{logo/logo_ubu.png}%
        }{%
            \fbox{%
                \parbox[c][2.8cm][c]{0.62\textwidth}{%
                    \centering
                    \textbf{LOGO UNIVERSIDAD DE BURGOS}\\[0.25cm]
                    Añadir \texttt{logo/logo\_ubu.png}%
                }%
            }%
        }

        \vspace{1.25cm}

        {\small\bfseries\color{UBUblue}\MakeUppercase{\grado}}\\[0.22cm]
        {\large\color{UBUdark}\universidad}

        \vspace{1.55cm}

        {\small\bfseries\color{UBUblue}\MakeUppercase{Asignatura}}\\[0.18cm]
        {\LARGE\bfseries\color{UBUdark}\asignatura}

        \vspace{0.75cm}
        {\color{UBUblue}\rule{0.22\textwidth}{3pt}}

        \vspace{1.0cm}

        {\small\bfseries\color{UBUblue}\MakeUppercase{Trabajo / práctica}}\\[0.30cm]
        {\Huge\bfseries\color{UBUdark}\titulotrabajo}

        \vfill

        % -------------------------------------------------
        % BLOQUE DE DATOS DEL ESTUDIANTE
        % -------------------------------------------------
        \noindent
        \colorbox{UBUlight}{%
            \parbox{0.96\linewidth}{%
                \vspace{0.35cm}
                \begin{tabularx}{\linewidth}{>{\leavevmode\bfseries\color{UBUdark}}p{4.0cm} X}
                    Estudiante & \nombrealumno \\[0.20cm]
                    DNI / NIE & \dni \\[0.20cm]
                    Correo electrónico & \correo \\[0.20cm]
                    Grupo & \grupo \\[0.20cm]
                    Profesor/a & \profesor \\[0.20cm]
                    Curso académico & \cursoacademico \\[0.20cm]
                    Fecha de entrega & \fechaentrega \\
                \end{tabularx}
                \vspace{0.15cm}
            }%
        }

        \vspace{0.85cm}

        {\small\color{UBUdark}
        Campus de Miranda de Ebro\hfill
        \textbf{Universidad de Burgos}}

        \vspace*{0.65cm}

    \end{minipage}
\end{titlepage}

% =========================================================
% ÍNDICE
% =========================================================

\tableofcontents
\newpage

% =========================================================
% 1. INTRODUCCIÓN
% =========================================================

\section{Introducción}

En esta práctica se parte de una simulación de una planta industrial desarrollada en
Python con \textbf{MuJoCo} (motor físico), \textbf{OpenCV} (visión artificial) y
\textbf{NumPy}. La planta está formada por cuatro cintas transportadoras que forman un
circuito cerrado, una quinta cinta que conduce las piezas hacia un horno, cámaras de
inspección situadas sobre cada cinta, varias piezas de distinta geometría y una isla
robotizada con un manipulador \textit{Franka Emika Panda} que dispone de cinemática
inversa (IK).

El código suministrado funcionaba, pero estaba escrito de forma procedural: un único
script con funciones, variables globales y estructuras de datos compartidas. El trabajo
realizado consiste en \textbf{refactorizar} ese código mediante Programación Orientada a
Objetos, de manera que cada elemento físico de la planta quede representado por una clase
con una responsabilidad clara, y que el resultado sirva como base para la Práctica~II.

Las funcionalidades principales de la aplicación final son:

\begin{itemize}[leftmargin=1.2cm]
    \item Transporte de piezas por un circuito de cuatro cintas con control individual de
          marcha y parada, más una cinta de salida hacia el horno.
    \item Inspección visual con cinco cámaras: captura de imagen, detección de la pieza y
          reconocimiento de su forma (círculo, cuadrado, rectángulo u otra) con OpenCV.
    \item Parada automática de la cinta cuando su cámara detecta una pieza, mediante la
          variable booleana \texttt{objeto\_detectado}.
    \item Horno con contador de piezas, control de temperatura y cambio de color en MuJoCo
          en función de la temperatura.
    \item Robot Panda que reutiliza la IK proporcionada sin modificarla.
    \item Ventana OpenCV con el mosaico de las cámaras, una vista general y una barra de
          estado, además del visor 3D de MuJoCo.
\end{itemize}

El lenguaje empleado es \textbf{Python 3}, con las bibliotecas \texttt{mujoco}~($\geq$3.2),
\texttt{numpy} y \texttt{opencv-python}. El desarrollo se ha realizado en PyCharm sobre
Windows.

% =========================================================
% 2. OBJETIVOS
% =========================================================
\clearpage
\section{Objetivos}

\subsection{Objetivo general}

Transformar la simulación procedural de la planta industrial en una arquitectura
orientada a objetos modular, mantenible y extensible, conservando todo su funcionamiento
(simulación en MuJoCo, visualización de cámaras con OpenCV y robot con IK) y añadiendo la
detección de formas, la parada automática de cintas y el control térmico del horno.

\subsection{Objetivos específicos}

\begin{enumerate}
    \item Identificar los elementos de la planta que deben representarse mediante clases.
    \item Diseñar atributos y métodos coherentes con la responsabilidad de cada clase.
    \item Eliminar las variables globales de estado (como \texttt{BANDERA\_CINTA\_1}) y
          encapsularlas dentro de los objetos.
    \item Utilizar composición entre objetos y, donde aporta claridad, herencia.
    \item Integrar el procesamiento de imagen con OpenCV dentro de la clase de cámara.
    \item Comunicar el resultado de la detección visual a la cinta correspondiente.
    \item Relacionar el estado interno del horno (temperatura) con su representación
          visual en MuJoCo.
    \item Validar el funcionamiento mediante las pruebas mínimas del enunciado.
\end{enumerate}

% =========================================================
% 3. ANÁLISIS DEL PROBLEMA
% =========================================================
\clearpage
\section{Análisis del problema}

\subsection{Descripción del problema}

La planta simula una línea de transporte en la que las piezas recorren el circuito
\textit{Cinta~1 $\rightarrow$ Cinta~2 $\rightarrow$ Cinta~3 $\rightarrow$ Cinta~4
$\rightarrow$ Cinta~1} en sentido horario. Desde la Cinta~4 las piezas pueden pasar a la
Cinta~5, que las conduce al horno. MuJoCo no dispone de una primitiva de cinta con
velocidad de superficie, por lo que el transporte se modela aplicando sobre cada pieza una
fuerza tangencial con realimentación de velocidad, una fuerza de centrado lateral y una
transición suave en las esquinas.

\subsection{Análisis del código procedural inicial}

El código de partida estaba organizado en un único fichero con funciones como
\texttt{build\_mjcf()}, \texttt{get\_ids()}, \texttt{reset\_objects()},
\texttt{update\_segment()}, \texttt{cinta\_activa()}, \texttt{conmutar\_cinta()},
\texttt{transport\_force()}, \texttt{apply\_transport()},
\texttt{render\_camera\_mosaic()} y \texttt{run()}. Al analizarlo se detectaron los
siguientes problemas:

\begin{itemize}
    \item \textbf{Estado global compartido}: la marcha de cada cinta se controlaba con
          variables globales \texttt{BANDERA\_CINTA\_1}, \texttt{BANDERA\_CINTA\_2}, \dots, y
          las funciones las consultaban o modificaban directamente.
    \item \textbf{Datos en estructuras paralelas}: la información de las piezas (ids de
          MuJoCo, segmento actual, altura de reposo) se guardaba en listas separadas
          indexadas por número de pieza, lo que obligaba a mantenerlas sincronizadas a mano.
    \item \textbf{Funciones con demasiadas responsabilidades}: \texttt{run()} construía el
          modelo, creaba ventanas, ejecutaba el bucle de simulación, renderizaba las
          cámaras y procesaba el teclado.
    \item \textbf{Lógica dispersa}: la física de cada cinta (dirección, centro, giro en la
          esquina) se decidía en \texttt{transport\_force()} mediante condicionales sobre el
          número de segmento, en lugar de residir en la propia cinta.
    \item \textbf{Ausencia de visión real}: las cámaras solo generaban imágenes para el
          mosaico; no se analizaban ni influían en la planta.
\end{itemize}

A partir de este análisis se asignó cada responsabilidad a un objeto: las funciones de
transporte pasan a \texttt{CintaTransportadora}, las de piezas a
\texttt{ObjetoIndustrial}, el render y la visión a \texttt{CamaraInspeccion}, y el bucle
\texttt{run()} se convierte en los métodos de \texttt{PlantaIndustrial}.

\begin{table}[H]
    \centering
    \caption{Correspondencia entre el código procedural y la versión orientada a objetos.}
    \label{tab:refactor}
    \begin{tabularx}{\textwidth}{p{4.6cm}X}
        \toprule
        \textbf{Código procedural} & \textbf{Versión POO} \\
        \midrule
        \texttt{BANDERA\_CINTA\_n} & Atributo \texttt{activa} de cada \texttt{CintaTransportadora} \\
        \texttt{cinta\_activa()}, \texttt{conmutar\_cinta()} & \texttt{esta\_activa()}, \texttt{arrancar()}, \texttt{detener()}, \texttt{conmutar()} \\
        \texttt{transport\_force()} & \texttt{CintaTransportadora.calcular\_fuerza()} \\
        \texttt{apply\_transport()} & \texttt{PlantaIndustrial.\_aplicar\_transporte()} + \texttt{ObjetoIndustrial.aplicar\_fuerza()} \\
        \texttt{update\_segment()} & \texttt{ObjetoIndustrial.actualizar\_cinta()} + \texttt{CintaTransportadora.ha\_terminado()} \\
        \texttt{reset\_objects()} & \texttt{PlantaIndustrial.reiniciar()} + \texttt{ObjetoIndustrial.colocar()} \\
        \texttt{get\_ids()} & Método \texttt{conectar(model)} de cada clase \\
        \texttt{build\_mjcf()} & \texttt{escena.construir\_mjcf()} + \texttt{generar\_xml()} de cada pieza \\
        \texttt{render\_camera\_mosaic()} & \texttt{CamaraInspeccion.capturar()/dibujar()} + \texttt{PlantaIndustrial.\_montar\_mosaico()} \\
        \texttt{run()} & \texttt{PlantaIndustrial.inicializar()}, \texttt{actualizar()}, \texttt{procesar\_teclado()}, \texttt{ejecutar()} \\
        \bottomrule
    \end{tabularx}
\end{table}

\subsection{Requisitos funcionales}

\begin{enumerate}
    \item La simulación debe ejecutarse en MuJoCo con visor 3D.
    \item Cada cinta debe tener su propio estado y poder arrancarse y detenerse de forma
          independiente (teclas \texttt{1}--\texttt{5}).
    \item Cada cámara debe capturar su imagen, mostrarla en OpenCV, detectar si hay una
          pieza y determinar su forma.
    \item La detección debe actualizar la variable booleana \texttt{objeto\_detectado} y,
          cuando sea verdadera, detener la cinta asociada.
    \item El horno debe registrar las piezas que entran, gestionar su temperatura y
          cambiar de color en MuJoCo según ésta.
    \item El robot Panda debe seguir funcionando con la IK proporcionada.
    \item Controles de teclado: pausa, reinicio, captura, encendido del horno, ajuste de
          temperatura y envío manual de una pieza al horno.
\end{enumerate}
\clearpage
\subsection{Restricciones y consideraciones}

\begin{itemize}
    \item No se acepta meter todo el código en una única clase: cada clase debe tener una
          responsabilidad clara.
    \item Debe reducirse al mínimo el uso de variables globales.
    \item La detección visual debe realizarse dentro de la clase de cámara.
    \item La cinemática inversa del robot \textbf{no} debe reimplementarse; se usa como
          caja negra.
    \item \texttt{main.py} debe quedar en la raíz y actuar solo como punto de entrada; la
          lógica va dentro de \texttt{src/}.
\end{itemize}

% =========================================================
% 4. DISEÑO ORIENTADO A OBJETOS
% =========================================================
\clearpage
\section{Diseño orientado a objetos}

El diseño sigue la idea de que \emph{cada elemento físico de la planta es un objeto que
conoce su propio estado y sabe actuar sobre MuJoCo}, mientras que
\texttt{PlantaIndustrial} se limita a crearlos, conectarlos y coordinarlos en cada paso de
la simulación.

\subsection{Identificación de clases}

\begin{table}[H]
    \centering
    \caption{Clases principales de la aplicación.}
    \label{tab:clases}
    \small
    \begin{tabularx}{\textwidth}{p{3.7cm}>{\raggedright\arraybackslash}X>{\raggedright\arraybackslash}X}
        \toprule
        \textbf{Clase} & \textbf{Responsabilidad} & \textbf{Elementos principales} \\
        \midrule
        \texttt{PlantaIndustrial} & Coordinar todos los elementos: construir el modelo, ejecutar el bucle principal, montar el mosaico y procesar el teclado. &
        \texttt{cintas}, \texttt{camaras}, \texttt{objetos}, \texttt{horno}, \texttt{robot}; \texttt{inicializar()}, \texttt{actualizar()}, \texttt{procesar\_teclado()}, \texttt{ejecutar()} \\
        \addlinespace
        \texttt{CintaTransportadora} & Guardar el estado y la geometría de una cinta y calcular la fuerza que ejerce sobre una pieza. &
        \texttt{numero}, \texttt{velocidad}, \texttt{direccion}, \texttt{activa}, \texttt{siguiente}; \texttt{arrancar()}, \texttt{detener()}, \texttt{calcular\_fuerza()} \\
        \addlinespace
        \texttt{CamaraInspeccion} & Capturar la imagen de una cámara MuJoCo, detectar la pieza y su forma con OpenCV y controlar su cinta. &
        \texttt{cinta}, \texttt{objeto\_detectado}, \texttt{forma\_detectada}; \texttt{capturar()}, \texttt{detectar\_forma()}, \texttt{actualizar()} \\
        \addlinespace
        \texttt{ObjetoIndustrial} & Representar una pieza: su geometría, su cuerpo en MuJoCo y la cinta por la que circula. &
        \texttt{forma}, \texttt{cinta}, \texttt{body\_id}; \texttt{generar\_xml()}, \texttt{colocar()}, \texttt{aplicar\_fuerza()} \\
        \addlinespace
        \texttt{Caja}, \texttt{Cilindro} & Subclases de \texttt{ObjetoIndustrial} con la geometría concreta. &
        \texttt{medidas} / \texttt{radio}; \texttt{generar\_xml()}, \texttt{estabilizar()} \\
        \addlinespace
        \texttt{Horno} & Registrar las piezas que entran, controlar la temperatura y reflejarla en el color del horno. &
        \texttt{temperatura}, \texttt{numero\_piezas}, \texttt{geom\_id}; \texttt{aumentar\_temperatura()}, \texttt{actualizar\_color()} \\
        \addlinespace
        \texttt{RobotPanda} & Envolver el módulo \texttt{robot.py} (IK proporcionada) para usarlo como un objeto más de la planta. &
        \texttt{config}, \texttt{estado}, \texttt{modo\_prueba\_ik}; \texttt{conectar()}, \texttt{controlar()}, \texttt{ir\_a\_objetivo()} \\
        \bottomrule
    \end{tabularx}
\end{table}

Además de las clases, el proyecto tiene tres módulos de apoyo que no modelan objetos del
dominio: \texttt{config.py} (todas las constantes de configuración), \texttt{escena.py}
(generación del XML de la escena) y \texttt{robot.py} (IK y máquina de estados del Panda).

\subsection{Diagrama de clases UML}

\begin{figure}[H]
    \centering
    \IfFileExists{figuras/diagrama_uml.png}{%
        \includegraphics[width=\textwidth]{diagrama_uml.png}%
    }{%
        \fbox{%
            \parbox[c][6cm][c]{0.82\textwidth}{%
                \centering
                Insertar aquí el diagrama de clases UML\\[0.3cm]
                \texttt{diagrama\_uml.png}%
            }%
        }%
    }
    \caption{Diagrama de clases UML de la aplicación.}
    \label{fig:uml}
\end{figure}

\subsection{Relaciones entre clases}

\begin{itemize}
    \item \textbf{Composición}: \texttt{PlantaIndustrial} crea y posee las cintas, las
          cámaras, las piezas, el horno y el robot. Su ciclo de vida depende de la planta
          (se crean en su constructor).
    \item \textbf{Asociación}: cada \texttt{CamaraInspeccion} guarda una referencia a la
          \texttt{CintaTransportadora} que vigila, y la controla llamando a sus métodos
          \texttt{detener()} y \texttt{arrancar()}. Cada \texttt{ObjetoIndustrial} guarda
          la cinta por la que circula en ese momento.
    \item \textbf{Autoasociación}: cada cinta apunta a la \texttt{siguiente}, formando el
          circuito $1\rightarrow2\rightarrow3\rightarrow4\rightarrow1$. Así una pieza pasa
          de cinta con \texttt{self.cinta = self.cinta.siguiente}, sin índices de segmento.
    \item \textbf{Herencia}: \texttt{Caja} y \texttt{Cilindro} heredan de
          \texttt{ObjetoIndustrial}.
    \item \textbf{Dependencia}: \texttt{RobotPanda} usa las funciones del módulo
          \texttt{robot.py}; \texttt{PlantaIndustrial} usa \texttt{escena.py} y todas las
          clases dependen de \texttt{config.py}.
\end{itemize}

% =========================================================
% 5. PRINCIPIOS DE POO APLICADOS
% =========================================================

\section{Principios de programación orientada a objetos aplicados}

\subsection{Encapsulación}

Cada objeto guarda su propio estado y lo expone mediante métodos. El ejemplo más claro es
la marcha de las cintas: las antiguas variables globales \texttt{BANDERA\_CINTA\_n} han
desaparecido y ahora cada cinta tiene su atributo \texttt{activa}, que se modifica con
\texttt{arrancar()}, \texttt{detener()} y \texttt{conmutar()}.

Los detalles internos se marcan con guion bajo (convención de Python para atributos
privados): por ejemplo, \texttt{\_ids\_rodillos} en la cinta, \texttt{\_qpos} y
\texttt{\_dof} en las piezas, o \texttt{\_inspeccionando} y \texttt{\_contorno} en la
cámara. También se usan \texttt{@property} para validar datos: la velocidad de una cinta
no puede ser negativa.

\begin{lstlisting}[caption={Velocidad de la cinta protegida mediante una propiedad.}]
@property
def velocidad(self):
    return self._velocidad

@velocidad.setter
def velocidad(self, valor):
    if valor < 0:
        raise ValueError("La velocidad de la cinta no puede ser negativa")
    self._velocidad = valor
\end{lstlisting}

En el horno, la temperatura solo cambia a través de \texttt{aumentar\_temperatura()} y
\texttt{disminuir\_temperatura()}, que además la mantienen entre la temperatura ambiente
y la máxima.

\subsection{Abstracción}

Las clases representan conceptos del dominio y ocultan cómo se implementan en MuJoCo. Desde
\texttt{PlantaIndustrial} basta con escribir \texttt{objeto.cinta.calcular\_fuerza(pos,
vel)} sin conocer el controlador interno, o \texttt{camara.actualizar(renderer, data)} sin
saber que por dentro hay una conversión a HSV, una apertura morfológica y un análisis de
contornos. Del mismo modo, \texttt{RobotPanda} oculta las funciones y el diccionario de
estado de \texttt{robot.py} detrás de métodos como \texttt{controlar()} o
\texttt{ir\_a\_home()}.

\subsection{Herencia}

Se ha utilizado una jerarquía sencilla para las piezas: \texttt{ObjetoIndustrial} contiene
todo lo común (conexión con MuJoCo, posición, fuerzas, cambio de cinta) y las subclases
\texttt{Caja} y \texttt{Cilindro} solo añaden lo que las distingue: sus medidas y su
geometría en el XML. Se eligió herencia porque la relación ``un cilindro \emph{es un}
objeto industrial'' es natural y evita duplicar código. En el resto de clases se ha
preferido la composición, ya que sus relaciones son de tipo ``\emph{tiene un}'' (la
planta tiene cintas, la cámara tiene una cinta asociada).

\subsection{Polimorfismo}

\texttt{generar\_xml()} se declara en la clase base lanzando
\texttt{NotImplementedError} (actúa como método abstracto) y cada subclase lo
sobrescribe. La planta construye la escena sin preguntar el tipo de cada pieza:

\begin{lstlisting}[caption={Uso polimórfico de \texttt{generar\_xml()}.}]
xml_objetos = "\n".join(obj.generar_xml() for obj in self.objetos)
\end{lstlisting}

También \texttt{estabilizar()} se sobrescribe en \texttt{Cilindro}: llama primero a la
versión de la clase base con \texttt{super()} y después centra el cilindro en la banda en
los tramos rectos, porque al rodar tiende a desviarse. Además, varias clases redefinen
\texttt{\_\_str\_\_()} para mostrar su estado en la barra de la ventana (por ejemplo,
\texttt{C1:ON} o \texttt{Horno 85C LISTO piezas=2}).

% =========================================================
% 6. IMPLEMENTACIÓN
% =========================================================
\clearpage
\section{Implementación}

\subsection{Estructura del proyecto}

\begin{itemize}
    \item \texttt{main.py}: punto de entrada; solo crea la planta y llama a
          \texttt{ejecutar()}.
    \item \texttt{requirements.txt}: dependencias (\texttt{mujoco}, \texttt{numpy},
          \texttt{opencv-python}).
    \item \texttt{src/PlantaIndustrial.py}: clase coordinadora y bucle principal.
    \item \texttt{src/CintaTransportadora.py}, \texttt{src/CamaraInspeccion.py},
          \texttt{src/ObjetoIndustrial.py}, \texttt{src/Horno.py},
          \texttt{src/RobotPanda.py}: una clase (o jerarquía) por fichero.
    \item \texttt{src/config.py}: constantes de configuración centralizadas.
    \item \texttt{src/escena.py}: generación del modelo MJCF de la planta.
    \item \texttt{src/robot.py}: IK y control del Panda (proporcionado).
    \item \texttt{src/vendor/mujoco\_menagerie/}: modelo del Franka Emika Panda, que se
          descarga automáticamente si no existe.
\end{itemize}

\begin{lstlisting}[caption={\texttt{main.py}: el programa principal queda separado de la lógica.}]
from src.PlantaIndustrial import PlantaIndustrial

def main():
    planta = PlantaIndustrial(mostrar_visor=True, modo_prueba_ik=True)
    planta.ejecutar()

if __name__ == "__main__":
    main()
\end{lstlisting}
\clearpage
\subsection{Clases y métodos principales}

\subsubsection{CintaTransportadora}

Cada cinta almacena su número, nombre, velocidad, dirección de transporte (vector
unitario), centro de la banda, dirección de giro en la esquina, referencia a la cinta
siguiente y su estado \texttt{activa}. Los parámetros del controlador (\texttt{VEL\_KP},
\texttt{CENTER\_KP}, \texttt{MAX\_FORCE}) son atributos de clase, ya que son iguales para
todas las cintas.

El método \texttt{calcular\_fuerza(pos, vel)} contiene la física que antes estaba en
\texttt{transport\_force()}: si la cinta está parada, frena la pieza; si está en marcha,
calcula la velocidad deseada (mezclando la dirección de la cinta y la de giro al acercarse
a la esquina), aplica un control proporcional sobre la velocidad y añade una fuerza de
centrado lateral. \texttt{actualizar\_rodillos()} hace girar o detiene los actuadores de
los rodillos según el estado, de modo que en el visor se ve qué cintas están paradas.

\subsubsection{ObjetoIndustrial, Caja y Cilindro}

Cada pieza guarda su número, nombre, color, masa, forma, la cinta actual y los índices de
MuJoCo (\texttt{body\_id}, \texttt{\_qpos}, \texttt{\_dof}) que obtiene en
\texttt{conectar(model)}. Ofrece métodos para leer su posición y velocidad, colocarla en un
punto, aplicarle fuerza, cambiar de cinta al llegar al final de la actual
(\texttt{actualizar\_cinta()}) y comprobar si se ha caído. Se simulan cinco piezas: un
cubo, dos cilindros y dos cajas rectangulares.
\clearpage
\subsubsection{PlantaIndustrial}

En el constructor crea todos los objetos a partir de las tablas de \texttt{config.py}
(\texttt{CINTAS}, \texttt{PIEZAS}, \texttt{NOMBRES\_CAMARAS}) y enlaza el circuito de
cintas. \texttt{inicializar()} construye el modelo MuJoCo, llama a \texttt{conectar()} en
cada objeto y abre las ventanas. Cada paso de simulación se ejecuta en
\texttt{actualizar()}:

\begin{lstlisting}[caption={Un paso de simulación en \texttt{PlantaIndustrial.actualizar()}.}]
def actualizar(self):
    if self.pausada:
        return
    self.robot.antes_del_paso(self.model, self.data, self.objetos, self.cintas)
    self._aplicar_transporte()
    self.robot.controlar(self.model, self.data)

    mujoco.mj_step(self.model, self.data)

    for objeto in self.objetos:
        objeto.estabilizar(self.data)
    self.robot.despues_del_paso(self.model, self.data)

    self.horno.actualizar(self.model.opt.timestep)
    self.horno.actualizar_color(self.model)
    for cinta in self.cintas:
        cinta.actualizar_rodillos(self.data)
    mujoco.mj_forward(self.model, self.data)
\end{lstlisting}

Las cámaras se actualizan a 20~fotogramas por segundo de tiempo simulado
(\texttt{actualizar\_camaras()}), y el bucle \texttt{ejecutar()} sincroniza el visor,
muestra el mosaico y procesa el teclado. La función \texttt{run()} original queda así
repartida en métodos cortos con una tarea cada uno.

\subsubsection{RobotPanda}

La IK no se ha modificado. \texttt{RobotPanda} actúa como \emph{adaptador}: guarda los
identificadores y el diccionario de estado que necesita \texttt{robot.py} y traduce los
objetos de la planta a las listas de números que esperan sus funciones (por ejemplo,
convierte \texttt{obj.cinta.numero} en el índice de segmento). Tiene dos modos: prueba de
IK (el robot va a un objetivo cartesiano, teclas \texttt{T} y \texttt{H}) y
\textit{pick-and-place} automático.

\subsection{Controles de teclado}

\begin{table}[H]
    \centering
    \caption{Teclas disponibles en la ventana de cámaras.}
    \label{tab:teclas}
    \begin{tabularx}{\textwidth}{p{2.6cm}X}
        \toprule
        \textbf{Tecla} & \textbf{Acción} \\
        \midrule
        \texttt{Q} / \texttt{ESC} & Salir. \\
        \texttt{P} & Pausar o reanudar la simulación. \\
        \texttt{1}--\texttt{5} & Conmutar la marcha de la cinta correspondiente. \\
        \texttt{R} & Reiniciar la planta (piezas, cintas, cámaras, horno y robot). \\
        \texttt{S} & Guardar una captura del mosaico en \texttt{capturas/}. \\
        \texttt{O} & Encender o apagar el horno. \\
        \texttt{+} / \texttt{-} & Subir o bajar la temperatura del horno 10\,\textdegree C. \\
        \texttt{E} & Enviar una pieza de la Cinta~4 a la Cinta~5 (horno). \\
        \texttt{T} / \texttt{H} & (Modo prueba IK) Mover el robot al objetivo / a la posición \textit{home}. \\
        \bottomrule
    \end{tabularx}
\end{table}

% =========================================================
% 6b. DETECCIÓN DE FORMAS Y PARADA DE CINTAS
% =========================================================
\clearpage
\section{Detección de formas y parada automática de las cintas}

\subsection{Mecanismo de detección de formas}

Toda la visión se realiza dentro de \texttt{CamaraInspeccion}. El método
\texttt{capturar()} renderiza la cámara MuJoCo asociada (\texttt{cam\_bottom},
\texttt{cam\_right}, \texttt{cam\_top}, \texttt{cam\_left} o \texttt{cam\_oven}) y
convierte la imagen de RGB a BGR para OpenCV. Después, \texttt{detectar\_forma()} aplica
estos pasos:

\begin{enumerate}
    \item \textbf{Segmentación por color} (\texttt{\_segmentar()}): la imagen se pasa a HSV.
          La banda, el suelo y las estructuras son grises (saturación baja), mientras que
          las piezas tienen colores vivos. Se seleccionan los píxeles con saturación
          $\geq 90$ y brillo $\geq 70$ con \texttt{cv2.inRange()}. Se eligió este método en
          lugar de una umbralización en escala de grises porque las sombras y los
          reflejos de la banda daban falsos positivos.
    \item \textbf{Limpieza morfológica}: una apertura con un núcleo de $5\times5$ elimina
          puntos sueltos de ruido.
    \item \textbf{Extracción de contornos} con \texttt{cv2.findContours()}.
    \item \textbf{Selección del contorno} (\texttt{\_elegir\_contorno()}): se descartan los
          contornos de área menor que 800 píxeles y los que tocan el borde de la imagen
          (la pieza aún no se ve entera). Se queda el de mayor área.
    \item \textbf{Clasificación} (\texttt{\_clasificar()}): se calcula la circularidad
          $C = 4\pi A / P^2$ (vale 1 en un círculo perfecto) y una aproximación poligonal
          con \texttt{cv2.approxPolyDP()}.
\end{enumerate}
\clearpage
\begin{lstlisting}[caption={Clasificación de la forma a partir del contorno.}]
def _clasificar(self, contorno):
    area = cv2.contourArea(contorno)
    perimetro = cv2.arcLength(contorno, True)
    circularidad = 4 * np.pi * area / (perimetro * perimetro)
    aproximacion = cv2.approxPolyDP(contorno, CAMARA_EPSILON_POLIGONO * perimetro, True)
    vertices = len(aproximacion)

    if circularidad > CAMARA_CIRCULARIDAD_MIN:
        return "circulo"
    # una caja vista un poco de lado enseña también una cara lateral,
    # por eso su contorno puede tener de 4 a 6 esquinas
    if 4 <= vertices <= 6:
        (_, _), (lado1, lado2), _ = cv2.minAreaRect(contorno)
        proporcion = max(lado1, lado2) / min(lado1, lado2)
        if proporcion < CAMARA_PROPORCION_CUADRADO:
            return "cuadrado"
        return "rectangulo"
    return "otra"
\end{lstlisting}

Así se distinguen \textbf{círculo} (cilindros vistos desde arriba, circularidad
$> 0{,}85$), \textbf{cuadrado} (cubo, proporción del rectángulo mínimo $< 1{,}25$),
\textbf{rectángulo} (cajas alargadas) y \textbf{otra} para cualquier contorno no
reconocido. Como las cámaras no están exactamente en la vertical, una caja muestra también
una cara lateral y su contorno puede tener 5 o 6 vértices; por eso se acepta el rango de 4
a 6. Todos los umbrales están en \texttt{config.py}, para poder ajustarlos sin tocar la
clase. El método \texttt{dibujar()} superpone el contorno detectado en verde y el texto
\textit{DETECTADO: forma}.

\subsection{Variable booleana \texttt{objeto\_detectado}}

Cada cámara tiene el atributo booleano \texttt{objeto\_detectado}, que se recalcula en cada
fotograma dentro de \texttt{actualizar()}: vale \texttt{True} cuando
\texttt{detectar\_forma()} devuelve una forma y \texttt{False} cuando devuelve
\texttt{None}.

\begin{lstlisting}[caption={Actualización de la cámara y comunicación con la cinta.}]
def actualizar(self, renderer, data):
    imagen = self.capturar(renderer, data)
    self.forma_detectada = self.detectar_forma(imagen)
    self.objeto_detectado = self.forma_detectada is not None
    self._controlar_cinta(data.time)
    return imagen
\end{lstlisting}
\clearpage
El método \texttt{\_controlar\_cinta()} usa esa variable para gobernar la cinta asociada.
Si se detendiera la cinta siempre que \texttt{objeto\_detectado} fuese verdadero, la pieza
quedaría parada para siempre bajo la cámara. Por eso se ha implementado un pequeño ciclo
de inspección controlado:

\begin{enumerate}
    \item Cuando aparece una pieza (\texttt{objeto\_detectado == True}) y no se estaba
          inspeccionando, se llama a \texttt{self.cinta.detener()}, se anota el instante y
          se incrementa \texttt{piezas\_inspeccionadas}.
    \item Mientras dura la inspección (1,5~s de tiempo simulado) la cinta permanece
          parada.
    \item Al terminar, la cámara llama a \texttt{self.cinta.arrancar()} y marca la pieza
          como \textit{liberada} (\texttt{\_pieza\_liberada}), para no volver a parar la
          cinta por la misma pieza mientras sale del campo de visión.
    \item Cuando la pieza sale de la imagen (\texttt{objeto\_detectado == False}) se borra
          la marca y la cámara queda lista para la siguiente pieza.
\end{enumerate}

De esta forma se cumple el comportamiento mínimo (\textit{si hay pieza, la cinta se
detiene}) y la decisión de reanudar la cinta queda aislada en un único método. En la
Práctica~II bastará con cambiar \texttt{\_controlar\_cinta()} (por ejemplo, para desviar
la pieza según su forma) sin modificar ni la visión ni la cinta.

% =========================================================
% 6c. HORNO
% =========================================================
\clearpage
\section{Clase Horno: temperatura y color en MuJoCo}

La clase \texttt{Horno} encapsula todo lo relativo al final de la línea. Sus atributos
principales son:

\begin{itemize}
    \item \texttt{numero\_piezas}: piezas recibidas, y \texttt{pieza\_nueva}: si acaba de
          entrar una pieza.
    \item \texttt{encendido}: estado del horno, y la propiedad \texttt{estado}, que devuelve
          \texttt{APAGADO}, \texttt{CALENTANDO} (por debajo de la temperatura mínima) o
          \texttt{LISTO}.
    \item \texttt{temperatura}: temperatura actual (inicial 20\,\textdegree C).
    \item \texttt{temperatura\_minima} (60\,\textdegree C) y \texttt{temperatura\_maxima}
          (220\,\textdegree C) de funcionamiento, y \texttt{temperatura\_alta}
          (150\,\textdegree C) para el cambio a rojo.
    \item \texttt{geom\_id}: identificador del \textit{geom} \texttt{oven\_roof} (techo del
          horno) en el modelo MuJoCo, que se obtiene en \texttt{conectar()}.
\end{itemize}

La temperatura solo se modifica mediante métodos:
\texttt{aumentar\_temperatura()} y \texttt{disminuir\_temperatura()} la limitan al rango
permitido, \texttt{actualizar(dt)} la hace subir 5\,\textdegree C/s si el horno está
encendido o bajar 4\,\textdegree C/s si está apagado, y \texttt{registrar\_pieza()} resta
5\,\textdegree C porque la pieza fría absorbe calor. Además, el usuario puede apagar o
encender el horno (\texttt{O}) y subir o bajar la temperatura manualmente (\texttt{+} y
\texttt{-}). Una pieza entra en el horno cuando, estando en la Cinta~5, su coordenada
$x$ supera la boca del horno (\texttt{detectar\_entrada()}); entonces se registra y la
pieza reaparece en la Cinta~1 para seguir circulando.

\begin{lstlisting}[caption={Cambio de color del horno según la temperatura.}]
def actualizar_color(self, model):
    if self.temperatura < self.temperatura_minima:
        color = self.AZUL
    elif self.temperatura < self.temperatura_alta:
        color = self.NARANJA
    else:
        color = self.ROJO
    model.geom_rgba[self.geom_id] = color
\end{lstlisting}

\begin{table}[H]
    \centering
    \caption{Codificación visual de la temperatura del horno.}
    \label{tab:horno}
    \begin{tabular}{lll}
        \toprule
        \textbf{Rango de temperatura} & \textbf{Estado} & \textbf{Color del techo} \\
        \midrule
        $T < 60$\,\textdegree C & CALENTANDO & Azul \\
        $60 \leq T < 150$\,\textdegree C & LISTO & Naranja \\
        $T \geq 150$\,\textdegree C & LISTO (alta) & Rojo \\
        \bottomrule
    \end{tabular}
\end{table}

\texttt{actualizar\_color()} se llama en cada paso de \texttt{PlantaIndustrial.actualizar()},
justo después de \texttt{horno.actualizar(dt)}, por lo que el cambio se ve en directo en el
visor de MuJoCo. Al modificar \texttt{model.geom\_rgba} se cambia el color del modelo, no
de un dato global, y todo el control térmico queda dentro de la clase.

% =========================================================
% 7. MANEJO DE ERRORES Y VALIDACIONES
% =========================================================

\section{Manejo de errores y validaciones}

\begin{itemize}
    \item \textbf{Validación de parámetros}: la velocidad de una cinta negativa lanza
          \texttt{ValueError}; un tipo de pieza desconocido en la configuración lanza
          \texttt{ValueError} al crear los objetos.
    \item \textbf{Elementos inexistentes en el modelo}: si no se encuentra un actuador de
          rodillo o el \textit{geom} del horno, \texttt{conectar()} lanza
          \texttt{RuntimeError} con el nombre que falta, en lugar de fallar más adelante
          con un índice $-1$.
    \item \textbf{Valores fuera de rango}: la temperatura del horno se limita entre la
          ambiente y la máxima; las fuerzas de transporte se recortan a
          \texttt{MAX\_FORCE}.
    \item \textbf{Piezas perdidas}: \texttt{se\_ha\_caido()} detecta si una pieza sale de
          los límites de la planta y la planta la recoloca en la Cinta~1.
    \item \textbf{Robot}: la máquina de estados proporcionada reinicia el ciclo si una
          fase de IK no converge en 7~s.
    \item \textbf{Modelo del Panda}: si no se encuentra \texttt{panda.xml}, se descarga
          automáticamente MuJoCo Menagerie; si falla, se muestra un mensaje explicando
          cómo clonarlo a mano.
    \item \textbf{Cierre ordenado}: \texttt{ejecutar()} usa \texttt{try/finally} para
          cerrar siempre el visor, el renderizador y las ventanas de OpenCV.
    \item \textbf{Acciones sin efecto}: si se pulsa \texttt{E} y no hay pieza en la
          Cinta~4, se informa por consola en lugar de producir un error.
\end{itemize}

% =========================================================
% 8. PRUEBAS
% =========================================================
\clearpage
\section{Pruebas y validación}

Se han realizado las pruebas mínimas que pide el enunciado ejecutando la simulación y
observando el visor, el mosaico de cámaras y los mensajes de consola.

\begin{table}[H]
    \centering
    \caption{Plan de pruebas.}
    \label{tab:pruebas}
    \small
    \begin{tabularx}{\textwidth}{p{1.2cm}>{\raggedright\arraybackslash}X>{\raggedright\arraybackslash}X p{1.5cm}}
        \toprule
        \textbf{Prueba} & \textbf{Entrada / escenario} & \textbf{Resultado esperado} & \textbf{Estado} \\
        \midrule
        P01 & Ejecutar \texttt{python main.py}. & Se abren el visor MuJoCo y la ventana de cámaras. & Correcto \\
        P02 & Pulsar \texttt{2} dos veces. & La Cinta~2 se detiene (rodillos parados, \texttt{C2:OFF}) y vuelve a arrancar. & Correcto \\
        P03 & Detener la Cinta~1 con el resto en marcha. & Solo se paran las piezas de la Cinta~1; las demás cintas siguen. & Correcto \\
        P04 & Observar el mosaico. & Las cinco cámaras y la vista general muestran imagen actualizada. & Correcto \\
        P05 & Pasar un cilindro bajo una cámara. & Se detecta \texttt{circulo}. & Correcto \\
        P06 & Pasar el cubo y las cajas bajo una cámara. & Se detecta \texttt{cuadrado} / \texttt{rectangulo}. & Correcto$^{*}$ \\
        P07 & Pieza entra en el campo de la cámara. & \texttt{objeto\_detectado} pasa a \texttt{True}; al salir, a \texttt{False}. & Correcto \\
        P08 & \texttt{objeto\_detectado == True}. & La cinta asociada se detiene y aparece \texttt{OFF} en la barra. & Correcto \\
        P09 & Fin de la inspección (1,5~s). & La cámara rearranca la cinta y no vuelve a pararla por la misma pieza. & Correcto \\
        P10 & Pulsar \texttt{+}, \texttt{-} y \texttt{O}. & La temperatura sube/baja 10\,\textdegree C; al apagar el horno se enfría. & Correcto \\
        P11 & Dejar calentar el horno. & El techo pasa de azul a naranja (60\,\textdegree C) y a rojo (150\,\textdegree C). & Correcto \\
        P12 & Pulsar \texttt{E}. & La pieza va a la Cinta~5, entra al horno, se cuenta y la temperatura baja 5\,\textdegree C. & Correcto \\
        P13 & Modo prueba IK, teclas \texttt{T} y \texttt{H}. & El Panda se mueve hacia el objetivo y hacia \textit{home} usando la IK proporcionada, sin reimplementarla. & \pendiente{Revisar} \\
        P14 & Pulsar \texttt{R}. & Piezas, cintas, cámaras, horno y robot vuelven al estado inicial. & Correcto \\
        \bottomrule
    \end{tabularx}
\end{table}

{\small $^{*}$ El cubo se clasifica como \texttt{cuadrado} en la mayoría de fotogramas,
pero en la Cinta~4, donde la cámara lo ve más inclinado, algún fotograma lo clasifica como
\texttt{rectangulo} porque la cara lateral alarga el contorno (ver
sección~\ref{sec:discusion}).}

% =========================================================
% 9. RESULTADOS
% =========================================================
\clearpage
\section{Resultados}

La Figura~\ref{fig:general} muestra la planta completa en el visor: el circuito de cuatro
cintas, la Cinta~5 hacia el horno, las cámaras sobre postes y el robot Panda en la isla
central.

\begin{figure}[H]
    \centering
    \includegraphics[width=0.92\textwidth]{overview_inicio.png}
    \caption{Vista general de la planta en MuJoCo.}
    \label{fig:general}
\end{figure}

La Figura~\ref{fig:mosaico} es la ventana de OpenCV durante la ejecución. Las cámaras de
las cintas 1 y 2 han detectado cilindros (\texttt{circulo}), la de la Cinta~3 una caja
(\texttt{rectangulo}) y la de la Cinta~4 el cubo (\texttt{cuadrado}). En la barra inferior
se ve que las cintas 1, 2 y 3 están en \texttt{OFF} porque sus cámaras están
inspeccionando (texto en rojo), mientras que la Cinta~4 ya ha terminado la inspección y
vuelve a estar en marcha (texto en verde).

\begin{figure}[H]
    \centering
    \includegraphics[width=\textwidth]{mosaico_deteccion.png}
    \caption{Mosaico de cámaras con detección de formas y parada automática de cintas.}
    \label{fig:mosaico}
\end{figure}

\begin{figure}[H]
    \centering
    \begin{subfigure}{0.32\textwidth}
        \centering
        \includegraphics[width=\textwidth]{det_circulo.png}
        \caption{Cilindro: círculo.}
    \end{subfigure}\hfill
    \begin{subfigure}{0.32\textwidth}
        \centering
        \includegraphics[width=\textwidth]{det_cuadrado.png}
        \caption{Cubo: cuadrado.}
    \end{subfigure}\hfill
    \begin{subfigure}{0.32\textwidth}
        \centering
        \includegraphics[width=\textwidth]{det_rectangulo.png}
        \caption{Caja: rectángulo.}
    \end{subfigure}
    \caption{Ejemplos de clasificación de las formas en las cámaras de inspección.}
    \label{fig:formas}
\end{figure}

La Figura~\ref{fig:horno} muestra el horno en el visor de MuJoCo a distintas temperaturas:
el techo (\texttt{oven\_roof}) cambia de azul a naranja y a rojo conforme el horno se calienta.

\begin{figure}[H]
    \centering
    \begin{subfigure}{0.32\textwidth}
        \centering
        \includegraphics[width=\textwidth]{horno_azul.png}
        \caption{$T < 60$\,\textdegree C: azul.}
    \end{subfigure}\hfill
    \begin{subfigure}{0.32\textwidth}
        \centering
        \includegraphics[width=\textwidth]{horno_naranja.png}
        \caption{$60 \leq T < 150$\,\textdegree C: naranja.}
    \end{subfigure}\hfill
    \begin{subfigure}{0.32\textwidth}
        \centering
        \includegraphics[width=\textwidth]{horno_rojo.png}
        \caption{$T \geq 150$\,\textdegree C: rojo.}
    \end{subfigure}
    \caption{Cambio de color del horno en MuJoCo según su temperatura.}
    \label{fig:horno}
\end{figure}

Por consola se registran los eventos de la planta, por ejemplo:

\begin{lstlisting}[language={},numbers=none,caption={Mensajes de consola durante la ejecución.}]
[CINTA 1 - INFERIOR] pieza detectada: circulo -> cinta 1 parada
[CINTA 4 - IZQUIERDA] pieza detectada: cuadrado -> cinta 4 parada
[CINTA 4 - IZQUIERDA] inspección terminada -> cinta 4 en marcha
obj0 (cubo) enviado a la Cinta 5
[Horno] entra obj0 (cubo) -> total 1
\end{lstlisting}

% =========================================================
% 10. DISCUSIÓN
% =========================================================

\section{Discusión}
\label{sec:discusion}

\textbf{Ventajas de la arquitectura.} Cada clase tiene una sola responsabilidad y se
comunica con las demás a través de métodos. La cámara no sabe cómo se mueve una cinta;
solo llama a \texttt{detener()} y \texttt{arrancar()}. La cinta no sabe que existe una
cámara. Gracias a esto, añadir un nuevo sensor o una nueva decisión en la Práctica~II no
obliga a cambiar el resto de clases.

\textbf{Configuración centralizada.} Todas las constantes (geometría, ganancias, umbrales
de visión, rangos de temperatura, piezas y posiciones iniciales) están en
\texttt{config.py}. Las piezas y las cintas se crean recorriendo tablas, de modo que
añadir una pieza nueva es añadir una línea a \texttt{PIEZAS}. Las únicas variables
globales que quedan son constantes, no estado de la planta.

\textbf{Encadenamiento de cintas.} Sustituir los índices de segmento por la referencia
\texttt{siguiente} eliminó los condicionales por número de cinta del código original:
cada cinta sabe por sí misma hacia dónde mueve las piezas y cuál le sigue.

\textbf{Adaptador del robot.} Envolver \texttt{robot.py} en \texttt{RobotPanda} permitió
cumplir el requisito de no reimplementar la IK y, a la vez, integrar el robot como un objeto
más. El inconveniente es que \texttt{robot.py} sigue trabajando con diccionarios y listas,
por lo que \texttt{RobotPanda} tiene que traducir los datos en cada paso.

\textbf{Decisiones que podrían mejorarse.}
\begin{itemize}
    \item La clasificación entre cuadrado y rectángulo depende de la perspectiva de la
          cámara. Para hacerla más robusta se podría corregir la perspectiva o usar varias
          imágenes seguidas y quedarse con la forma más repetida.
    \item \texttt{PlantaIndustrial} sigue teniendo bastantes tareas (bucle, mosaico,
          teclado). Se podrían separar en clases como \texttt{SistemaVision} (mosaico) o
          \texttt{ControladorPlanta} (teclado y decisiones).
    \item La comunicación cámara--cinta es directa. Para más sensores, un controlador
          intermedio que reciba las detecciones y decida qué hacer reduciría el
          acoplamiento.
    \item Python no impide acceder a atributos con guion bajo; la encapsulación se basa en
          la convención.
\end{itemize}

\textbf{Extensiones futuras (Práctica~II).} Clasificar las piezas según su forma o color
y enviarlas al horno solo si cumplen una condición, que el robot recoja únicamente
determinadas piezas, o que el horno solo acepte piezas cuando su estado sea
\texttt{LISTO}. Todas estas decisiones pueden añadirse en \texttt{\_controlar\_cinta()} o
en un nuevo controlador sin tocar la física ni la visión.

% =========================================================
% 11. CONCLUSIONES
% =========================================================

\section{Conclusiones}

Se ha transformado la simulación procedural de la planta en una arquitectura orientada a
objetos con las clases \texttt{PlantaIndustrial}, \texttt{CintaTransportadora},
\texttt{CamaraInspeccion}, \texttt{ObjetoIndustrial} (con \texttt{Caja} y
\texttt{Cilindro}), \texttt{Horno} y \texttt{RobotPanda}. La simulación sigue
funcionando en MuJoCo con la visualización de cámaras en OpenCV y el robot con la IK
proporcionada.

Se han cumplido los objetivos planteados: las banderas globales de las cintas han
desaparecido y cada cinta gestiona su estado; las cámaras detectan las piezas y
reconocen su forma con OpenCV; la variable \texttt{objeto\_detectado} detiene la cinta
correspondiente y la reanudación se decide de forma controlada; y el horno gestiona su
temperatura con métodos propios y la muestra mediante el color en MuJoCo.

Más allá de que el programa funcione, la práctica ha servido para comprobar en un caso
real cómo la encapsulación, la composición y la separación de responsabilidades hacen que
el código sea más fácil de entender, probar y ampliar. El resultado deja preparada una
base modular para la Práctica~II.

% =========================================================
% 12. CONTRIBUCIÓN INDIVIDUAL
% =========================================================

\section{Contribución individual}

\begin{table}[H]
    \centering
    \caption{Contribución de los integrantes del grupo.}
    \begin{tabularx}{\textwidth}{>{\raggedright\arraybackslash}p{4cm}X}
        \toprule
        \textbf{Estudiante} & \textbf{Contribución} \\
        \midrule
        \nombrealumno & Análisis del código procedural, diseño de clases, refactorización, implementación de la detección de formas y del control del horno, pruebas y redacción del informe. \\
        \bottomrule
    \end{tabularx}
\end{table}

% =========================================================
% 13. USO DE INTELIGENCIA ARTIFICIAL
% =========================================================
\clearpage
\section{Declaración sobre el uso de herramientas de inteligencia artificial}

\begin{quote}
Para hacer este informe me he apoyado en una IA generativa (Claude), sobre todo para la redacción, la maquetación en LaTeX y para sacar el diagrama de clases. Aun así, he revisado todo lo que me ha generado y lo he corregido donde hacía falta, para asegurarme de que el contenido está bien y refleja el trabajo que he hecho.
\end{quote}

% =========================================================
% 14. REFERENCIAS
% =========================================================

\addcontentsline{toc}{section}{Referencias}
\begin{thebibliography}{9}

\bibitem{enunciado}
Universidad de Burgos,
\textit{Práctica I: Refactorización de una Planta Industrial mediante Programación
Orientada a Objetos}, Programación Modular y Orientada a Objetos, curso 2026--2027.

\bibitem{mujoco}
E. Todorov, T. Erez e Y. Tassa,
``MuJoCo: A physics engine for model-based control'',
\textit{IEEE/RSJ International Conference on Intelligent Robots and Systems}, 2012.
\url{https://mujoco.readthedocs.io}

\bibitem{menagerie}
Google DeepMind,
\textit{MuJoCo Menagerie: Franka Emika Panda}.
\url{https://github.com/google-deepmind/mujoco_menagerie}

\bibitem{opencv}
OpenCV,
\textit{Contour Features}, documentación de OpenCV.
\url{https://docs.opencv.org}

\bibitem{gamma}
E. Gamma, R. Helm, R. Johnson y J. Vlissides,
\textit{Design Patterns: Elements of Reusable Object-Oriented Software},
Addison-Wesley.

\bibitem{martin}
R. C. Martin,
\textit{Clean Code: A Handbook of Agile Software Craftsmanship},
Prentice Hall.

\bibitem{uml}
Object Management Group,
\textit{Unified Modeling Language (UML)}.

\end{thebibliography}

% =========================================================
% ANEXOS
% =========================================================

\newpage
\appendix

\section{Código fuente y repositorio}

El código completo del proyecto está disponible en el repositorio GitHub:

\begin{center}
    \large\url{\repositorio}
\end{center}

\textbf{Instalación y ejecución:}

\begin{lstlisting}[language=bash,numbers=none]
python -m venv .venv
.venv\Scripts\activate          # en Linux/macOS: source .venv/bin/activate
pip install -r requirements.txt
python main.py
\end{lstlisting}

La primera vez, si no se encuentra el modelo del Panda, el programa descarga
automáticamente la carpeta \texttt{franka\_emika\_panda} de MuJoCo Menagerie dentro de
\texttt{src/vendor/} (requiere \texttt{git}). También se puede indicar una copia local con
la variable de entorno \texttt{MUJOCO\_MENAGERIE\_PATH}.

Para probar el \textit{pick-and-place} automático del robot en lugar del modo de prueba de
la IK, basta con crear la planta con \texttt{PlantaIndustrial(modo\_prueba\_ik=False)} en
\texttt{main.py}.

\section{Material adicional}

\subsection*{Parámetros de visión}

\begin{table}[H]
    \centering
    \caption{Parámetros de detección definidos en \texttt{config.py}.}
    \small
    \begin{tabularx}{\textwidth}{p{6.6cm}lX}
        \toprule
        \textbf{Constante} & \textbf{Valor} & \textbf{Uso} \\
        \midrule
        \texttt{CAMARA\_UMBRAL\_SATURACION} & 90 & Saturación mínima (HSV) de una pieza \\
        \texttt{CAMARA\_UMBRAL\_BRILLO} & 70 & Brillo mínimo (HSV) de una pieza \\
        \texttt{CAMARA\_AREA\_MINIMA} & 800 px & Área mínima del contorno \\
        \texttt{CAMARA\_MARGEN\_BORDE} & 3 px & Descarta contornos que tocan el borde \\
        \texttt{CAMARA\_NUCLEO\_MORFOLOGIA} & 5 & Tamaño del núcleo de apertura \\
        \texttt{CAMARA\_CIRCULARIDAD\_MIN} & 0,85 & Umbral para \texttt{circulo} \\
        \texttt{CAMARA\_EPSILON\_POLIGONO} & 0,04 & Tolerancia de \texttt{approxPolyDP} \\
        \texttt{CAMARA\_PROPORCION\_CUADRADO} & 1,25 & Límite cuadrado / rectángulo \\
        \texttt{CAMARA\_TIEMPO\_INSPECCION} & 1,5 s & Tiempo de parada por pieza \\
        \bottomrule
    \end{tabularx}
\end{table}

\end{document}
